\documentclass[
  oneside,
  author={Shijie Lu},
  degree={BSc},
  title={Label-Efficient Semantic Segmentation via Self-Supervised Depth Estimation and Cross-Task Consistency},
  unit={COMS30045},
  subtitle={}
]{dissertation}
\usepackage{booktabs}
\usepackage{amssymb}
\usepackage{tabularx}
\usepackage{makecell}
\usepackage{float}
\usepackage{units}
\usepackage{url}
\begin{document}

\maketitle

\frontmatter

% -----------------------------------------------------------------------------
\chapter*{Abstract}

\noindent
Semantic segmentation requires dense pixel-level annotation, averaging
around 1.5~hours per image in complex urban scenes, making full dataset
annotation prohibitively expensive. This motivates research into training
strong segmentation models under limited annotation budgets.
Hoyer et al.\ proposed a framework that integrates self-supervised depth
estimation~(SDE) into semi-supervised segmentation through three
complementary components: automatic data selection, geometry-aware mixing
augmentation~(DepthMix), and multi-task learning, achieving strong
performance under label-efficient conditions.

This project hypothesises that introducing an explicit cross-task
feature alignment objective~($\mathcal{L}_\text{ct}$) between the depth
and segmentation decoder branches of the PAD-Net multi-task architecture
can produce stable, reproducible improvements under a constrained
single-GPU academic computing environment, and that this benefit remains
additive when combined with the existing framework components.

The main contributions of this project are as follows:

\begin{itemize}

\item The four result tables of \cite{hoyer2023ijcv_sde} (data
  selection, mixing strategy, SDE feature transfer, and framework
  component ablation) were successfully reproduced on Cityscapes under
  a single-GPU constraint at the University of Bristol. All qualitative
  performance trends are preserved, with most configurations falling
  within $\pm$1\,mIoU of the paper.

\item $\mathcal{L}_\text{ct}$ was designed and implemented within the
  PAD-Net multi-task decoder, with two loss formulations---normalised
  MSE and cosine similarity---evaluated through a systematic two-phase
  calibration procedure.

\item The best-performing configuration (single-projection cosine loss,
  $\lambda_\text{ct}=1.00$) yields a consistent $+2.50$\,pp improvement
  over the PAD-MTL baseline across three independent seeds, with a
  standard deviation of only $0.17$\,\%.

\item When combined with the full framework, the strongest configuration
  achieves $68.19$\,\% mIoU at the 372-label regime, exceeding both the
  reproduced full-framework baseline by $+0.53$\,pp and the original
  paper's multi-GPU result by $+0.18$\,pp, demonstrating that the
  proposed method produces genuine improvements even under more
  constrained hardware conditions.

\end{itemize}

% -----------------------------------------------------------------------------
\chapter*{Dedication and Acknowledgements}

\vspace{1cm}

\noindent
I would like to thank my supervisor, Weihong Li, for his support and
advice as I completed this project.

% -----------------------------------------------------------------------------
\makedecl

% -----------------------------------------------------------------------------
\makeaidecl

% -----------------------------------------------------------------------------
\tableofcontents
\listoffigures
\listoftables

% -----------------------------------------------------------------------------
\chapter*{Ethics Statement}

This project did not require ethical review, as determined by my
supervisor, Weihong Li.

% -----------------------------------------------------------------------------
\chapter*{Supporting Technologies}

\noindent
The following third-party resources were used during the execution and
write-up of this project.

\begin{itemize}

    \item \textbf{Isambard-AI:} The University of Bristol's
    Isambard-AI high-performance computing platform was used to run the
    reproduction experiments (exp~210--213), corresponding to the data
    selection, mixing strategy, feature transfer, and component ablation
    results. Each job was allocated a single GPU and managed via the
    Slurm workload manager.

    \item \textbf{BluePebble (BP1):} The University of Bristol's
    BluePebble HPC cluster was used to run all cross-task consistency
    loss experiments (exp~219--222). Jobs were submitted via Slurm and
    executed predominantly on NVIDIA V100 nodes with 32\,GB VRAM.
    Each job was allocated one GPU, eight CPU cores, and 64\,GB of
    system memory.

    \item \textbf{Slurm Workload Manager:} Slurm was used on both HPC
    platforms to submit, schedule, and monitor all training and
    evaluation jobs throughout the project.

    \item \textbf{Cityscapes Dataset:} The Cityscapes dataset provided
    all training and evaluation data. Three components were used:
    \texttt{leftImg8bit} (RGB images), \texttt{gtFine} (pixel-level
    semantic annotations), and \texttt{leftImg8bit\_sequence} (unlabelled
    video sequences for self-supervised depth estimation).

    \item \textbf{PyTorch:} The PyTorch deep learning framework was used
    to implement, train, and evaluate all neural network models, including
    the segmentation decoder, depth decoder, and the proposed cross-task
    consistency loss module.

    \item \textbf{Python Scientific Stack:} NumPy was used for numerical
    computation and array manipulation. Matplotlib was used to produce
    training curves and result visualisations. Additional Python utilities
    handled data loading, logging, and experiment configuration via a
    centralised \texttt{experiments.py} module.

    \item \textbf{Conda:} A dedicated Conda environment with pinned
    package versions was used on both HPC platforms to ensure consistent
    software dependencies across all node types and job submissions.

    \item \textbf{Git:} Git was used for version control throughout the
    project, enabling reproducible tracking of code changes across all
    experiment configurations and SLURM submission scripts.

\end{itemize}

% -----------------------------------------------------------------------------
\chapter*{Notation and Acronyms}

\noindent
\begin{tabular}{lcl}
HPC         &:& High Performance Computing                        \\
BP1         &:& BluePebble Phase~1 (University of Bristol HPC)   \\
GPU         &:& Graphics Processing Unit                          \\
CPU         &:& Central Processing Unit                           \\
VRAM        &:& Video Random-Access Memory                        \\
SSL         &:& Semi-Supervised Learning                          \\
SSS         &:& Semi-Supervised Semantic Segmentation             \\
SSDA        &:& Semi-Supervised Domain Adaptation                 \\
SDE         &:& Self-Supervised Monocular Depth Estimation        \\
MTL         &:& Multi-Task Learning                               \\
mIoU        &:& Mean Intersection-over-Union                      \\
IoU         &:& Intersection-over-Union                           \\
CNN         &:& Convolutional Neural Network                      \\
ASPP        &:& Atrous Spatial Pyramid Pooling                    \\
EMA         &:& Exponential Moving Average                        \\
PL          &:& Pseudo-Label(s)                                   \\
CE          &:& Cross-Entropy loss                                \\
MSE         &:& Mean Squared Error                                \\
S           &:& Data Selection (framework component)              \\
DX          &:& DepthMix (framework component)                    \\
DS          &:& Diversity Sampling                                \\
US          &:& Uncertainty Sampling                              \\
DA          &:& Data Augmentation                                 \\
RGB         &:& Red, Green, Blue image channels                   \\
GT          &:& Ground Truth                                      \\
fd0         &:& Standard SDE encoder pre-training variant         \\
fd2         &:& Feature-distance regularised SDE pre-training variant \\
lr          &:& Learning rate                                     \\
$\lambda$   &:& Loss weighting coefficient                        \\
$\mathcal{L}_\text{ct}$ &:& Cross-Task Consistency Loss           \\
$\theta_T$  &:& Teacher network weights                           \\
$\theta_S$  &:& Student network weights                           \\
$\alpha$    &:& EMA smoothing coefficient                         \\
\end{tabular}

% =============================================================================
\mainmatter

\include{chapter1}
\include{chapter2}
\include{chapter3}
\include{chapter4}

% -----------------------------------------------------------------------------
\chapter{Conclusion}
\label{chap:conclusion}

% ============================================================
\section{Summary of Contributions}

\noindent
This project had two primary goals: to reproduce the SDE-enhanced
label-efficient segmentation framework of \cite{hoyer2023ijcv_sde} under
a constrained single-GPU environment, and to extend the PAD-Net
multi-task architecture with a Cross-Task Consistency Loss
($\mathcal{L}_\text{ct}$) that provides an explicit feature-alignment
objective between the segmentation and depth decoder branches. All
experiments were conducted under a single-GPU constraint with batch
size~2 and without Synchronised Batch Normalisation, which introduces
a systematic downward bias relative to the authors' multi-GPU setup.
Within these constraints, the reproduction successfully recovers the
qualitative trends of the original paper, and $\mathcal{L}_\text{ct}$
produces consistent, seed-stable improvements that remain positive when
combined with the full framework.

\begin{itemize}

\item \textbf{Faithful reproduction of four result tables.}
  The data selection (Table~1), mixing strategy (Table~3), SDE feature
  transfer (Table~5), and framework component ablation (Table~7) from
  \cite{hoyer2023ijcv_sde} were reproduced on Cityscapes. The data
  selection experiment covered multiple label regimes, while the main
  Stage~(b) comparisons were reproduced at the 1/8-label (372-label)
  regime. All qualitative performance trends are preserved, with most
  configurations falling within $\pm$1\,mIoU of the paper.

\item \textbf{Design and implementation of $\mathcal{L}_\text{ct}$.}
  A lightweight cross-task projection module was implemented within the
  PAD-Net multi-task decoder. Two loss formulations---normalised MSE and
  cosine similarity---were evaluated through a two-phase calibration
  procedure, incurring only $16{,}512$ additional parameters with no
  extra forward pass through the depth branch.

\item \textbf{Calibration and multi-seed validation.}
  The best stable configuration (single $d{\rightarrow}s$ cosine,
  $\lambda_\text{ct}=1.00$) achieves $63.25$\,\% mIoU versus a PAD-MTL
  baseline of $60.75$\,\%, corresponding to a $+2.50$\,pp improvement
  that is stable and reproducible across three independent seeds
  (std~$0.17$\,\%), with a stable operating range of
  $\lambda_\text{ct} \in [0.75, 1.00]$.

\item \textbf{Full framework composability (exp~222).}
  $\mathcal{L}_\text{ct}$ was evaluated on MTL\,+\,DX, MTL\,+\,S, and
  MTL\,+\,DX\,+\,S. All six configurations showed positive mean-level
  gains. The strongest configuration achieved $68.19$\,\% mIoU,
  exceeding the reproduced full-framework baseline by $+0.53$\,pp and
  the paper's multi-GPU result by $+0.18$\,pp.

\end{itemize}

% ============================================================
\section{Evaluation Against Aims and Objectives}

\noindent
Chapter~\ref{chap:context} set out four specific objectives, all of
which were completed.

\begin{enumerate}

\item \textbf{Reproduce Stage~(a)~[Completed].}
  All four data selection strategies (Random, Entropy, DS, DS+US) were
  verified, with DS+US consistently achieving the highest mIoU.
  Absolute values are within $1.2$\,pp of the paper across all label
  regimes, and the trend that informed selection improves over random
  sampling is consistently reproduced.

\item \textbf{Reproduce Stage~(b)~[Completed].}
  Tables~3, 5, and~7 were reproduced, confirming the qualitative trends
  for mixing strategy, SDE transfer, and framework component ablation.
  The most notable deviation is in MTL\,+\,DX (Table~7), where the
  reproduced mean falls $1.55$\,pp below the paper value, likely related
  to the reduced batch size, which may degrade batch statistics,
  pseudo-label stability, and the implicit cross-task signal that
  DepthMix relies on.

\item \textbf{Design, implement, and evaluate
  $\mathcal{L}_\text{ct}$~[Completed].}
  The proposed loss was calibrated over two phases and validated across
  three independent seeds (Section~\ref{sec:ct-calibration-ablation}).
  The best stable configuration achieves $63.25$\,\% mIoU versus a
  PAD-MTL baseline of $60.75$\,\%, corresponding to a $+2.50$\,pp
  improvement that is stable and reproducible across seeds.

\item \textbf{Evaluate the full combined framework~[Completed].}
  Exp~222 confirmed that $\mathcal{L}_\text{ct}$ composes positively
  with all three component combinations. The complete system achieved
  $68.19$\,\% mIoU, the highest result in this project and comparable
  to the paper's reported $68.01$\,\%.

\end{enumerate}

% ============================================================
\section{Limitations and Future Work}

\noindent
Three directions emerge naturally from the experimental findings.

\paragraph{Component-conditioned $\lambda$ scheduling.}
The composability ablation showed that the preferred loss type is
component-dependent: within the $d{\rightarrow}s$ track, MSE achieves
higher gains when DepthMix is absent (MTL\,+\,S: $+0.60$\,pp), while
cosine is more stable when DepthMix is present (MTL\,+\,DX: $+0.97$\,pp).
The MSE variant ultimately achieves the highest raw mean on
MTL\,+\,DX\,+\,S ($68.19$\,\%), but with higher variance
(std~$1.03$\,\%) than the cosine counterpart. Treating
$\lambda_\text{ct}$ as a component-conditioned hyperparameter---or
annealing it when DepthMix is active---could reduce this variance
without sacrificing the mean-level gain.

\paragraph{Dual projection revisited.}
Dual projection was only evaluated at the same $\lambda_\text{ct}$
values used for single projection. Because dual projection introduces a
symmetric alignment constraint and potentially larger alignment
gradients, lower values (e.g., $\lambda \in \{0.25, 0.50\}$) may yield
a more stable training regime. A targeted sweep in this region could
determine whether dual projection is fundamentally inferior or merely
over-regularised at the current $\lambda$ range.

\paragraph{Generalisation beyond Cityscapes.}
The proposed $\mathcal{L}_\text{ct}$ was evaluated only on Cityscapes
at the 1/8-label regime, where the structural correlation between depth
and segmentation is particularly strong due to the urban driving domain.
Whether the loss transfers to domains with weaker geometric
structure---such as indoor scenes or aerial imagery---remains open.
Evaluating on datasets such as ADE20K~\cite{zhou2019ade20k}, PASCAL
VOC, or indoor RGB-D benchmarks, provided that suitable self-supervised
or pretrained depth estimates are available, would provide a direct test
of domain generality. A longer-term direction is to extend the same
alignment idea to additional geometrically related auxiliary tasks, such
as surface normal estimation or optical flow.

% =============================================================================
\bibliography{dissertation}

% -----------------------------------------------------------------------------
\appendix

\chapter{Appendix A: AI Prompts/Tools}
\label{appx:ai_prompt}

\noindent
The following list describes all uses of AI tools during this project
and its write-up.

\begin{itemize}

    \item \textbf{Claude (Anthropic) --- \LaTeX{} Formatting Assistance:}
    Claude was used to help format \LaTeX{} tables, align equations, and
    structure front-matter sections including the Notation and Acronyms
    table and the Supporting Technologies page. An example prompt:
    \textit{``The following is a draft LaTeX table for my reproduction
    results. Please fix the column alignment and spacing so that the
    reproduced and reported values are vertically aligned within each
    cell.''}

    \item \textbf{Claude (Anthropic) --- SLURM Script Debugging:}
    Claude was used to debug issues in SLURM job submission scripts, such
    as resolving environment activation errors, correcting path references,
    and adjusting wall-time and memory allocation settings. An example
    prompt: \textit{``Here is the error log from a failed SLURM job on
    BluePebble. The job exits immediately after the conda activate step.
    What is the likely cause and how should I fix the submission
    script?''}

    \item \textbf{Claude (Anthropic) --- Result Figure Generation:}
    Claude was used to generate Python Matplotlib code for producing the
    result visualisation figures in Chapter~\ref{chap:evaluation},
    including the Phase~II three-seed validation chart and the
    component-level ablation chart. Experimental result data were
    provided alongside the desired figure title and axis labels, and
    Claude generated the corresponding plotting script. An example
    prompt: \textit{``Here are the Phase~II three-seed validation
    results. Please generate a horizontal bar chart sorted by mean
    mIoU, with the selected configuration highlighted.$\pm$
    Std by Configuration'.''}

    \item \textbf{Grammarly / Overleaf Grammar Checker:}
    The built-in Overleaf grammar and spell checker was used during
    the write-up to identify typographic errors and grammatical issues
    in the dissertation text.

\end{itemize}

\chapter{Code Repository}
\label{appx:code}

\noindent
The full source code for this project, including all experiment
configurations, training scripts, cross-task consistency loss
implementation, and SLURM job submission scripts, is available at:

\begin{center}
\url{https://github.com/ShijieLu222/SDE_SEG.git}
\end{center}

% =============================================================================
\end{document}